% ============================================================
% MỤC 1: GIỚI THIỆU (đúng thứ tự khung của cô)
%   tổng quan - giải thuật liên quan - hiện trạng thị trường - lý do chọn - đối tượng/đầu ra - mục tiêu - phạm vi - tính năng - timeline
% Số liệu thị trường lấy từ trang giới thiệu của từng sản phẩm (xem mục Tài liệu tham khảo [9]-[11]).
% Chỗ cần ảnh dùng \slideimg: chưa có ảnh thì hiện khung xám ghi tên file.
% ============================================================
\section{Giới Thiệu}

% ---------- Tổng quan ----------
\begin{frame}[label=gt]{Giới Thiệu Đề Tài}
  \begin{columns}[T]
    \begin{column}{0.54\textwidth}
      \begin{block}{Battleship (Bắn tàu) nhiều người chơi qua mạng}
        \begin{itemize}
          \item Hai người chơi, mỗi người một lưới $10\times10$ và một đội tàu.
          \item Hai bên \textbf{lần lượt} gọi tọa độ để bắn; mỗi phát bắn là \textbf{trúng} hoặc \textbf{hụt}.
          \item Ai đánh chìm hết tàu đối phương trước thì thắng.
          \item Vị trí tàu của đối phương luôn bị \textbf{che} (sương mù).
        \end{itemize}
      \end{block}
      \begin{exampleblock}{Vì sao chọn đề tài này cho NT106}
        Trò chơi theo lượt, dữ liệu nhỏ, nhưng đủ các bài toán mạng thật: kết nối, đồng bộ trạng thái, chống gian lận, ngắt kết nối giữa chừng.
      \end{exampleblock}
    \end{column}
    \begin{column}{0.42\textwidth}
      \centering
      % [ẢNH] hình minh họa luật chơi / màn hình trận đấu
      \slideimg[\textwidth]{gk-01-man-hinh-tran.png}{Màn hình bàn cờ 10x10 của game (chụp ở Client). Nếu chưa có gameplay thì chụp lưới 10x10.}
    \end{column}
  \end{columns}
\end{frame}

% ---------- Giải thuật (1/2): mạng và đồng thời ----------
\begin{frame}{Giải Thuật Liên Quan (1/2): Mạng Và Đồng Thời}
  \footnotesize
  \begin{tabularx}{\textwidth}{@{}l l Y@{}}
    \toprule \textbf{Giải thuật} & \textbf{Dùng ở} & \textbf{Ý tưởng chính} \\ \midrule
    Đọc khung theo độ dài (length-prefix) & \texttt{PacketFraming} &
      Đọc đủ 4 byte độ dài $N$, rồi \textbf{lặp đọc} tới khi đủ $N$ byte mới giải mã. Giải quyết dính gói, vỡ gói của TCP. \\
    Hàng đợi FIFO ghép cặp & \texttt{MatchmakingQueue} &
      \texttt{ConcurrentQueue} + một worker nền: cứ có $\ge 2$ người thì lấy \textbf{hai người đầu hàng} ghép thành phòng, ai vào trước được ghép trước. Kiểm tra còn kết nối ngay trước khi ghép. \\
    Hủy tìm trận kiểu ``đánh dấu'' & \texttt{MatchmakingQueue} &
      Hàng đợi không xóa được phần tử ở giữa, nên khi hủy chỉ \textbf{xóa tên khỏi bảng băm} những người còn chờ; lúc lấy ra mà không thấy tên thì bỏ qua. Hủy mất $O(1)$. \\
    Định tuyến gói theo kiểu & \texttt{PacketRouter} &
      Tra bảng băm \texttt{Type} $\rightarrow$ \texttt{Handler}, mất $O(1)$. Thêm tính năng chỉ cần đăng ký thêm một dòng. \\
    Mỗi Client một Task & \texttt{TcpServer} &
      \texttt{async/await}: chờ mạng không chặn luồng; \texttt{SemaphoreSlim} bảo đảm mỗi Client chỉ có một Task ghi vào stream. \\
    \bottomrule
  \end{tabularx}
\end{frame}

% ---------- Giải thuật (2/2): bảo mật và luật chơi ----------
\begin{frame}{Giải Thuật Liên Quan (2/2): Bảo Mật Và Luật Chơi}
  \footnotesize
  \begin{tabularx}{\textwidth}{@{}l l Y@{}}
    \toprule \textbf{Giải thuật} & \textbf{Dùng ở} & \textbf{Ý tưởng chính} \\ \midrule
    AES-128-CBC & \texttt{CryptoHelper} &
      Mã hóa khối đối xứng, \textbf{IV ngẫu nhiên mỗi gói}; khóa 16 byte lấy từ SHA-256 của mật khẩu chung. \\
    BCrypt & \texttt{AuthHandler} &
      Băm một chiều có \textbf{salt} và hệ số chi phí 11 (chuỗi lưu dạng \texttt{\$2a\$11\$...}), cố tình chậm để chống dò mật khẩu. \\
    Giới hạn tần suất (cửa sổ cố định) & \texttt{LoginRateLimiter} &
      Đếm số lần sai \textbf{theo IP và theo tên}; đủ 5 lần thì khóa 60 giây. Đăng ký tối đa 5 lần mỗi 10 phút mỗi IP. \\
    Xử lý lượt bắn (thiết kế) & Turn Engine &
      Kiểm tra đúng lượt và ô chưa bắn, tra bản đồ theo ô trong $O(1)$; đếm số ô tàu còn lại để biết chìm hay thắng. \textbf{Trúng thì giữ lượt, hụt thì đổi lượt}; quá 30 giây tự đổi lượt. \\
    Hàng gửi riêng từng client (thiết kế) & ClientSession &
      Mỗi client có hàng đợi gửi và một Task ghi riêng; phát gói cho nhiều người chỉ cần bỏ vào hàng, người chậm không làm chậm người khác. \\
    \bottomrule
  \end{tabularx}
\end{frame}

% ---------- Hiện trạng thị trường ----------
\begin{frame}[label=hientrang]{Hiện Trạng: Các Game Battleship Trực Tuyến Có Sẵn}
  \scriptsize
  \begin{tabularx}{\textwidth}{@{}l Y Y Y Y@{}}
    \toprule
    \textbf{Tiêu chí} & \textbf{Battleship (battleship-game.org) \cite{ref9}} & \textbf{Hasbro's BATTLESHIP (Steam) \cite{ref10}} & \textbf{Fleet Battle (di động) \cite{ref11}} & \textbf{Đồ án của nhóm} \\
    \midrule
    Nền tảng & Web, chạy trên trình duyệt & PC (Steam) & Android, iOS & Windows (WinForms) \\
    Chơi qua mạng & 2 người; chọn bạn bè hoặc đối thủ ngẫu nhiên & Online PvP, chế độ Classic và Commanders & Ghép trận nhanh với người chơi khắp nơi, có giải đấu & Tạo/vào phòng bằng ID, tìm trận ngẫu nhiên \\
    Chat & Có, chat với đối thủ & Không thấy nêu & Có, chat chung và theo ngôn ngữ & Chat sảnh (chat trong trận: đang làm) \\
    Bảng xếp hạng & Không thấy nêu & Không thấy nêu & Có & Có (REST Supabase) \\
    Chi phí & Chạy ngay trên web & Trả phí (khoảng 2,5 -- 5 USD tùy nơi bán) & Có quảng cáo và mua trong ứng dụng & Miễn phí \\
    Tầng mạng & Chưa thấy tài liệu công bố & Chưa thấy tài liệu công bố & Chưa thấy tài liệu công bố & \textbf{Tự xây trên TCP thô, có tài liệu giao thức} \\
    \bottomrule
  \end{tabularx}
  \vspace{0.15cm}

  \tiny Thông tin lấy từ trang giới thiệu của từng sản phẩm (truy cập 09/11/2026), nhóm chưa tự cài đặt để kiểm chứng; ``không thấy nêu'' nghĩa là trang giới thiệu không đề cập.
\end{frame}

% ---------- Lý do chọn đề tài ----------
\begin{frame}{Cái Đã Có, Cái Chưa Có Và Lý Do Chọn Đề Tài}
  \begin{columns}[T]
    \begin{column}{0.48\textwidth}
      \begin{block}{Các sản phẩm trên đã làm được}
        \begin{itemize}
          \item Chơi \textbf{trực tuyến thật} với người lạ hoặc bạn bè.
          \item Có chat, bảng xếp hạng, giải đấu (tùy sản phẩm).
          \item Chạy ổn định ở quy mô lớn.
        \end{itemize}
      \end{block}
      \begin{alertblock}{Chưa phục vụ mục tiêu học tập mạng}
        \begin{itemize}
          \item \textbf{Không xem được tầng mạng}: giao thức, cách đóng gói, cách chống gian lận.
          \item Một số sản phẩm \textbf{trả phí} hoặc có \textbf{quảng cáo}.
          \item Không thể mở rộng hay thử nghiệm tính năng mới.
        \end{itemize}
      \end{alertblock}
    \end{column}
    \begin{column}{0.48\textwidth}
      \begin{exampleblock}{Vì sao nhóm tự xây}
        \begin{itemize}
          \item Môn NT106 yêu cầu làm việc \textbf{trực tiếp với socket}, nên cần tự cài giao thức.
          \item Battleship theo lượt: luật đơn giản, nhưng đủ các vấn đề mạng thật (đồng bộ trạng thái, chống gian lận, rớt kết nối).
          \item Có chỗ để thêm điểm riêng: \textbf{tối ưu hiệu năng và chịu tải}, \textbf{CSDL trên cloud}.
        \end{itemize}
      \end{exampleblock}
    \end{column}
  \end{columns}
\end{frame}

% ---------- Đối tượng + đầu ra ----------
\begin{frame}{Đối Tượng Áp Dụng Và Đầu Ra Mong Đợi}
  \begin{columns}[T]
    \begin{column}{0.46\textwidth}
      \begin{block}{Đối tượng áp dụng}
        \begin{itemize}
          \item \textbf{Người chơi} muốn chơi Battleship qua mạng.
          \item \textbf{Người học lập trình mạng}: dùng làm ví dụ về socket, giao thức, bảo mật.
          \item \textbf{Người vận hành} (nhóm): chạy Server, quản lý dữ liệu.
        \end{itemize}
      \end{block}
      \begin{exampleblock}{Vai trò trong nhóm}
        \begin{tabular}{@{}ll@{}}
          Đức & Server, phòng, luật trận \\
          Dũng & Giao diện, cơ sở dữ liệu \\
          Hoàng & Giao thức, mạng, bảo mật
        \end{tabular}
      \end{exampleblock}
    \end{column}
    \begin{column}{0.5\textwidth}
      \begin{block}{Đầu ra}
        \begin{enumerate}
          \item \textbf{Server} TCP bất đồng bộ (C\# .NET 8).
          \item \textbf{Client} WinForms chơi được từ đăng nhập đến hết trận.
          \item \textbf{Cơ sở dữ liệu} trên cloud (PostgreSQL) kèm REST đọc bảng xếp hạng.
          \item \textbf{Giao thức mạng} có tài liệu: khung gói tin, bắt tay, mã hóa.
          \item Báo cáo, slide và bản demo.
        \end{enumerate}
      \end{block}
    \end{column}
  \end{columns}
\end{frame}

% ---------- Mục tiêu ----------
\begin{frame}[label=mt]{Mục Tiêu Đề Tài}
  \begin{block}{Mục tiêu chung}
    Xây dựng game Battleship hai người chơi qua mạng, \textbf{tự cài giao thức trên TCP thô}, an toàn và chạy được giữa hai mạng khác nhau.
  \end{block}
  \begin{columns}[T]
    \begin{column}{0.48\textwidth}
      \begin{exampleblock}{Mục tiêu cụ thể (kỹ thuật)}
        \begin{enumerate}\setlength{\itemsep}{0pt}
          \item Giao thức riêng: khung gói 4 byte, bắt tay, định tuyến.
          \item Server xử lý nhiều người chơi đồng thời, không sập khi Client rớt mạng.
          \item \textbf{Server nắm luật}, chống gian lận.
          \item Bảo mật: AES-128, BCrypt, khóa dò mật khẩu.
        \end{enumerate}
      \end{exampleblock}
    \end{column}
    \begin{column}{0.48\textwidth}
      \begin{exampleblock}{Mục tiêu cụ thể (sản phẩm)}
        \begin{enumerate}\setlength{\itemsep}{0pt}
          \item Đăng ký, đăng nhập, phòng, ghép trận, chat.
          \item Chơi trọn một trận: xếp tàu, bắn theo lượt, thắng thua.
          \item Lưu lịch sử và bảng xếp hạng trên cloud.
          \item Tính năng riêng: tối ưu hiệu năng và chịu tải nhiều người chơi.
        \end{enumerate}
      \end{exampleblock}
    \end{column}
  \end{columns}
\end{frame}

% ---------- Phạm vi ----------
\begin{frame}{Phạm Vi Đề Tài}
  \begin{columns}[T]
    \begin{column}{0.5\textwidth}
      \begin{block}{Trong phạm vi}
        \begin{itemize}
          \item Tài khoản: đăng ký, đăng nhập, trạng thái Online/Offline.
          \item Sảnh chờ: chat, danh sách phòng, tạo/vào phòng bằng ID.
          \item Ghép trận tự động (Matchmaking).
          \item Xếp tàu, bắn theo lượt, sương mù, xử thắng khi đối thủ mất kết nối.
          \item Lưu lịch sử trận, bảng xếp hạng.
          \item Bảo mật: AES-128, BCrypt, khóa đăng nhập.
          \item Tính năng đặc biệt: \textbf{tối ưu hiệu năng và chịu tải nhiều người chơi}.
        \end{itemize}
      \end{block}
    \end{column}
    \begin{column}{0.46\textwidth}
      \begin{alertblock}{Ngoài phạm vi}
        \begin{itemize}
          \item Chế độ khán giả (Spectator).
          \item Đối thủ máy (AI), bản di động, đồ họa 3D.
          \item Thanh toán, hệ thống bạn bè nâng cao.
        \end{itemize}
      \end{alertblock}
      \begin{exampleblock}{Nền tảng}
        Ứng dụng desktop \textbf{Windows} (WinForms, C\# .NET 8). Chạy trên máy cá nhân; demo giữa hai mạng khác nhau bằng Ngrok.
      \end{exampleblock}
    \end{column}
  \end{columns}
\end{frame}

% ---------- Chức năng Server (có dấu tích) ----------
\newcommand{\tickDone}{\textcolor{green!50!black}{$\checkmark$}}
\newcommand{\tickHalf}{\textcolor{orange!80!black}{$\circ$}}
\newcommand{\tickNone}{\textcolor{gray}{$-$}}

\begin{frame}[label=tinhnangsv]{Chức Năng Phía Server}
  \scriptsize
  \renewcommand{\arraystretch}{1.05}
  \begin{tabularx}{\textwidth}{@{}l Y c@{}}
    \toprule \textbf{Nhóm} & \textbf{Chức năng} & \textbf{Trạng thái} \\ \midrule
    Mạng & Nhận kết nối TCP bất đồng bộ, mỗi Client một Task & \tickDone \\
    Mạng & Khung gói 4 byte, bắt tay kiểm tra phiên bản giao thức & \tickDone \\
    Mạng & Mã hóa AES-128-CBC mọi gói tin & \tickDone \\
    Mạng & PacketRouter chuyển gói tới handler của từng chức năng & \tickDone \\
    Tài khoản & Đăng ký, đăng nhập, băm mật khẩu BCrypt & \tickDone \\
    Tài khoản & Quên mật khẩu: gửi mã xác nhận qua email (SMTP), đặt lại mật khẩu & \tickNone \\
    Tài khoản & Đăng xuất: dọn phiên, gỡ khỏi Online, hủy hàng đợi, rời phòng & \tickNone \\
    Bảo mật & Khóa 60 giây sau 5 lần sai, giới hạn số lần đăng ký & \tickDone \\
    Bảo mật & Chặn gói khi chưa đăng nhập, một tài khoản một phiên & \tickDone \\
    Sảnh & Danh sách Online/Offline, chat sảnh & \tickDone \\
    Phòng & Tạo phòng, vào phòng, rời phòng & \tickDone \\
    Phòng & Kick thành viên (có hàm, chưa nối vào Router) & \tickHalf \\
    Ghép trận & Hàng đợi ghép cặp tự động & \tickDone \\
    Dữ liệu & SQLite / PostgreSQL (Supabase), view bảng xếp hạng & \tickDone \\
    Dữ liệu & Ghi lịch sử trận (có bảng, chưa ghi dữ liệu) & \tickHalf \\
    Trận đấu & Ready, xếp tàu, bắn theo lượt, GameOver, timer 30 giây, xử thắng khi rớt mạng & \tickNone \\
    Đặc biệt & Tối ưu hiệu năng, chịu tải nhiều người chơi (hàng gửi riêng, mã hóa một lần, đo trước và sau) & \tickNone \\
    \bottomrule
  \end{tabularx}
  \par\vspace{0.05cm}
  {\tiny \tickDone{} xong \quad \tickHalf{} làm dở (đã có mã, chưa dùng được đầy đủ) \quad \tickNone{} chưa làm (theo lịch ở slide Timeline)}
\end{frame}

% ---------- Chức năng Client (có dấu tích) ----------
\begin{frame}[label=tinhnangcl]{Chức Năng Phía Client}
  \scriptsize
  \renewcommand{\arraystretch}{1.05}
  \begin{tabularx}{\textwidth}{@{}l Y c@{}}
    \toprule \textbf{Nhóm} & \textbf{Chức năng} & \textbf{Trạng thái} \\ \midrule
    Mạng & NetworkManager: kết nối, bắt tay, vòng nhận gói chạy nền & \tickDone \\
    Mạng & Mã hóa AES mọi gói gửi và nhận & \tickDone \\
    Tài khoản & Màn hình đăng nhập, khóa có đếm ngược 60 giây & \tickDone \\
    Tài khoản & Màn hình đăng ký & \tickDone \\
    Tài khoản & Quên mật khẩu: xin mã, đặt lại; thêm ô email khi đăng ký & \tickNone \\
    Tài khoản & Nút đăng xuất, quay về màn hình đăng nhập & \tickNone \\
    Tài khoản & Hồ sơ người chơi (Profile) & \tickDone \\
    Sảnh & Lobby và chat sảnh & \tickDone \\
    Phòng & Tạo phòng, danh sách phòng, vào phòng bằng ID & \tickDone \\
    Ghép trận & Nút Find Opponent, vào phòng chờ khi được ghép & \tickDone \\
    Phòng chờ & Hiện tên hai người chơi & \tickDone \\
    Phòng chờ & Nút Ready đồng bộ với đối thủ (hiện chỉ đổi chữ) & \tickNone \\
    Thống kê & Bảng xếp hạng đọc qua REST & \tickDone \\
    Trận đấu & Bàn cờ $10\times10$ (giao diện) & \tickDone \\
    Trận đấu & Đặt, xoay tàu, kiểm tra hợp lệ (chưa nối mạng) & \tickHalf \\
    Trận đấu & Bắn, hiện trúng/hụt, sương mù, chat trong trận & \tickNone \\
    \bottomrule
  \end{tabularx}
  \par\vspace{0.05cm}
  {\tiny \tickDone{} xong \quad \tickHalf{} làm dở (đã có mã, chưa dùng được đầy đủ) \quad \tickNone{} chưa làm (theo lịch ở slide Timeline)}
\end{frame}

% ---------- Timeline ----------
\begin{frame}[label=timeline]{Timeline Cập Nhật Đến Hết Học Kỳ}
  \centering
  \begin{tikzpicture}[
    ph/.style={draw=inkblue, rounded corners=3pt, thick, align=center, font=\scriptsize,
               minimum height=1.9cm, text width=2.35cm, inner sep=3pt},
    arr/.style={-{Stealth[length=2mm]}, thick, color=brandblue}]
    \node[ph, fill=paleblue] (a) at (0,0)    {\textbf{10/09 -- 01/11}\\Nền tảng, mạng, đăng nhập, phòng, DB cloud\\[1pt]\textcolor{green!50!black}{\textbf{Phần lớn đã xong}}};
    \node[ph, fill=gold!40]  (b) at (2.75,0) {\textbf{02/11 -- 15/11}\\Tuần thi (02--08/11), báo cáo giữa kỳ, Matchmaking};
    \node[ph, fill=white]    (c) at (5.5,0)  {\textbf{16/11 -- 02/12}\\Gameplay chia nhiều nhịp: xếp tàu, luật trận, tối ưu hiệu năng, quên mật khẩu, sương mù\\[1pt]\textbf{Đóng băng code 02/12}};
    \node[ph, fill=white]    (d) at (8.25,0) {\textbf{03/12 -- 08/12}\\Kiểm thử, Ngrok, báo cáo Chương 4--5};
    \node[ph, fill=white]    (e) at (11.0,0) {\textbf{09/12 -- 14/12}\\Đóng gói, chạy thử, hoàn thiện báo cáo};
    \draw[arr] (a) -- (b); \draw[arr] (b) -- (c); \draw[arr] (c) -- (d); \draw[arr] (d) -- (e);
    \node[font=\scriptsize, align=center] at (5.5,-1.8) {Sẵn sàng báo cáo từ \textbf{15/12/2026}; buổi học cuối 25/12/2026};
  \end{tikzpicture}

  \vspace{0.3cm}
  \begin{alertblock}{Ghi chú}
    \footnotesize Việc then chốt là \textbf{logic trận đấu trên Server} (nhận Ready, xếp tàu, bắn, kết thúc). Các mục sương mù và lịch sử trận đều phụ thuộc vào mục này.
  \end{alertblock}
\end{frame}
